% ECONOMICS_PROBLEM: {"id": "Q54-512", "title": "Cost-effective climate adaptation", "jel_code": "Q54", "source_id": "OP-0433"}
% !TeX program = xelatex
\documentclass[11pt,a4paper]{article}
\usepackage[margin=23mm]{geometry}
\usepackage{fontspec}
\setmainfont{Latin Modern Roman}
\usepackage{amsmath,amssymb,mathtools}
\usepackage{xcolor,enumitem,booktabs,longtable,array,xurl,hyperref}
\definecolor{muted}{HTML}{53636C}
\hypersetup{colorlinks=true,urlcolor=blue,linkcolor=blue}
\setlength{\parindent}{0pt}
\setlength{\parskip}{5pt}
\newcommand{\R}{\mathbb R}
\newcommand{\E}{\mathbb E}
\newcommand{\N}{\mathbb N}
\newcommand{\ind}{\mathbf 1}
\newcommand{\Var}{\operatorname{Var}}
\newcommand{\Cov}{\operatorname{Cov}}
\newcommand{\supp}{\operatorname{supp}}
\DeclareMathOperator*{\argmin}{arg\,min}
\DeclareMathOperator*{\argmax}{arg\,max}
\newcommand{\entrymeta}[3]{{\sffamily\small\color{muted}\textbf{\texttt{#1}}\quad|\quad #2\par #3\par}}
\newcommand{\Problemref}[1]{\texttt{#1}}
\newcommand{\JELref}[2]{\texttt{#2} (JEL \texttt{#1})}
\begin{document}
\section*{Cost-effective climate adaptation}
\textbf{Problem ID:} \texttt{Q54-512}\quad \textbf{JEL:} \texttt{Q54}\par
% BEGIN ECONOMICS STATEMENT
\label{problem:OP-0433}
{\sffamily\small\color{muted}Original subject group: Climate, environment and energy\par}
\label{definitions:problem:30007715}
\textbf{Audit classification:} Published research agenda. Checked source identity and appropriate agenda/background support; expanded intervention, units, population, definedness and causal restrictions. Exact targets and variants are editorial.
\subsection*{Definitions and precise research target}
\textit{Editorial formalization.} Consider low-income households in heat-exposed districts of one country and a thirty-year sequence of projected temperatures. Let \(a\in\mathcal A\) denote an implementable adaptation package, where the finite set \(\mathcal A\) includes cooling access, warning systems, and infrastructure upgrades. Define household discounted health and consumption utility \(U_i(a,c)\) under climate path \(c\), and discounted public resource cost \(K(a,c)\). The climate distribution \(P\) and welfare weights \(w_i>0\) must be stated before estimation. Include a zero-cost baseline package, assume finite expected utilities and costs, and let \(B\geq0\) be the budget. The target is
\[a^*(P,B)\in\arg\max_{a\in\mathcal A:E_P[K(a,c)]\leq B}E_P\!\left[\sum_i w_iU_i(a,c)\right].\]
Assume packages can be assigned at district level, but allow household responses and within-district spillovers. Estimate policy effects using credible intervention variation and identify which response parameters must be transported beyond observed weather. The open research agenda concerns welfare under future climates, rather than merely whether present-day cooling reduces heat harm. Compare rankings across plausible climate distributions, uptake costs, and later complementary investments. If available variation does not identify long-run responses, give explicit bounds or alternative scenario rankings. All package definitions, budgets, utility normalization, and transport restrictions are part of this proposed operationalization.

Let \(I\) be the finite baseline household cohort and \(c=(T_0,\ldots,T_{30},\xi)\) the climate and economic shock path under law \(P\). Define \(U_i(a,c)=\sum_{t=0}^{30}\beta^tu_i(c_{it}^{\rm cons}(a,c),h_{it}(a,c))\), with consumption and health units and utility normalizations declared. \(K\) is real incremental public cost discounted with stated resource prices; sunk baseline services are common to all alternatives. Require \(K(0,c)=0\), integrable outcomes, and positive feasible consumption. The finite feasible menu is nonempty, so an optimizer exists when mean responses are known. The substantive question is causal recovery and transport of those responses and their welfare rankings, not unresolved existence of a finite argmax. All expectations use a stated baseline population and probability law. Identification means every admissible data-generating model with the same observed-data law yields the same displayed target; otherwise report the identified set. Exact equations, horizons and assumptions are editorial, rather than source-given canonical conjectures.
\subsection*{Variants and assumption changes}
\begin{enumerate}
\item \textbf{Editorial robust-policy variant.} Maximize \(\inf_{P\in\mathcal P}E_P\sum_iw_iU_i(a,c)\), subject to \(\sup_{P\in\mathcal P}E_PK(a,c)\leq B\), over a declared climate-law ambiguity set with a feasible baseline.
\item \textbf{Editorial dynamic-package variant.} Replace one fixed package with nonanticipating \(a_t\) contingent on observed heat and uptake. Compare adaptive and fixed policies under the same resource budget and information.
\end{enumerate}
\subsection*{References and source audit}
\textbf{1.} Explicit policy-oriented adaptation research agenda; exact interventions and optimizers are editorial. Locator: VoxEU column, 15 April 2025; introductory summary. Primary indexed author summary checked; subsequent full-text retrieval failed, so precise subsection claims are not independently reaffirmed. URL: \url{https://cepr.org/voxeu/columns/economics-climate-adaptation-academic-insights-effective-policy}.
\begin{enumerate}\raggedright
\item\hypertarget{definitions:source:30007715:1}{} Tamma Carleton; Esther Duflo; Kelsey Jack; Guglielmo Zappalà. 2025. \emph{The economics of climate adaptation: From academic insights to effective policy}. VoxEU column, 15 April 2025; introductory summary.
\url{https://cepr.org/voxeu/columns/economics-climate-adaptation-academic-insights-effective-policy}.
\end{enumerate}

\subsection*{Common conventions: Source mathematical conventions}

These conventions apply to each entry unless explicitly replaced.

\textbf{Models and identification.} A primitive model \(m\) specifies all probability laws, feasible choices, information, equilibrium and measurement rules used by its target. The admissible class \(\mathcal M\) is fixed independently of outcomes. If \(P_O(m)\) is its observable probability law and \(\tau(m)\) the target, its identified set at observable law \(P\) is
\[
 \mathcal I_\tau(P)=\{\tau(m):m\in\mathcal M,\ P_O(m)=P\}.
\]
Point identification means that this set is a singleton. An empty set signals incompatibility of data and assumptions. Coordinate bounds need not describe the joint set. Bounds are sharp only if every retained target is attainable or arbitrarily approachable by compatible models. Finite bounds require explicit restrictions on unobserved quantities; unrestricted confounding, hidden wealth, extrapolation or arbitrary equilibrium selection can yield uninformative sets.

\textbf{Causal interventions.} A potential outcome \(Y_i(p)\) is the outcome under a complete policy regime \(p\), including timing, financing, information and market spillovers. The target population and units are fixed before treatment. With interference, use the complete assignment vector or a justified exposure mapping; individual no-interference notation alone cannot identify an economy-wide intervention. A difference of mean potential outcomes does not require the cross-world joint distribution. Conditioning on post-treatment migration, survival, enrollment or employment changes the estimand and needs separate justification.

\textbf{Statistical statements.} An estimator, confidence set, forecast or test specifies a sampling law, model class and loss/coverage criterion. Uniform \((1-\alpha)\)-coverage means \(\inf_{m\in\mathcal M}\Pr_m\{\tau(m)\in C_n\}\geq1-\alpha\), or an explicitly asymptotic version. The whole parameter space trivially has coverage, so informativeness requires a stated width, risk or power objective. A forecasting improvement is relative to a named benchmark, information set and evaluation population.

\textbf{Equilibrium and optimization.} An equilibrium comprises feasible plans, optimality, beliefs where needed, budget constraints and all market/resource clearing. The primitives or a stated benchmark define each correspondence \(\mathcal E(p;m)\). Nonemptiness is assumed or proved, never inferred just from compact policy sets. A selected-equilibrium target uses a declared selection rule; a robust target quantifies over all permitted equilibria. Use suprema or infima unless attainment is established. An \(\varepsilon\)-optimal policy is feasible and within \(\varepsilon\) of the finite optimum value. Report an empty feasible set separately.

\textbf{Welfare and accounting.} Normative utility, interpersonal normalization, social weights, numeraire, discounting and the use of public receipts are primitives. Behavioral utility need not coincide with the welfare criterion. Household welfare includes ownership income and rents once; adding profits again to compensating variation generally double-counts them. A monetary equivalent is defined by an explicit transfer and price convention, with monotonicity and range conditions. Average effects, mechanism contributions, transfers and welfare losses are different targets. Infinite sums require convergence and dynamic feasible plans require terminal or no-Ponzi conditions.

\textbf{Source versions.} A working-paper date, revision date and journal publication year are distinguished. A DOI for a journal version is not automatically the DOI of an earlier manuscript. Locators refer to the stated version and printed pages where specified. A metadata-only check does not verify a claim about a full-text conclusion. Every entry's source subsection narrows the provenance claim accordingly.
% END ECONOMICS STATEMENT
\end{document}
